%% =====================================================================
%%  Multimodal Stroke Prediction System
%%  IEEE Conference Paper  --  compile with pdfLaTeX (Overleaf: IEEEtran)
%% =====================================================================
\documentclass[conference]{IEEEtran}
\IEEEoverridecommandlockouts

\usepackage{cite}
\usepackage{amsmath,amssymb,amsfonts}
% IEEE requires Times. Under a Unicode engine (XeTeX/LuaTeX) the legacy `times`
% package leaves TU/ptm undefined and silently substitutes Latin Modern, so use the
% newtx family, which provides Times-compatible text and maths on all engines.
% newtxmath must follow amssymb, and \Bbbk is defined by both.
% T1 encoding is required for the ogonek and other accents in the bibliography,
% and is the encoding newtx expects.
\usepackage[T1]{fontenc}
\let\Bbbk\relax
\usepackage{newtxtext,newtxmath}
\usepackage{algorithmic}
\usepackage{graphicx}
\usepackage{textcomp}
\usepackage{xcolor}
\usepackage{booktabs}
\usepackage{url}
\def\BibTeX{{\rm B\kern-.05em{\sc i\kern-.025em b}\kern-.08em
    T\kern-.1667em\lower.7ex\hbox{E}\kern-.125emX}}

\begin{document}

\title{A Multimodal Stroke Prediction System Using\\
Smartphone Photoplethysmography, Self-Reported Risk Factors,\\
and Medical Document Analysis}

%% =====================================================================
%%  AUTHOR BLOCK -- fill in team member details here
%% =====================================================================
\author{
\IEEEauthorblockN{
    \rule{3cm}{0.4pt}\quad
    \rule{3cm}{0.4pt}\quad
    \rule{3cm}{0.4pt}
}
\IEEEauthorblockA{
    \textit{\rule{4cm}{0.4pt}} \\
    \textit{\rule{5cm}{0.4pt}} \\
    \rule{4cm}{0.4pt} \\
    \rule{5cm}{0.4pt}
}
\vspace{6mm}
}
%% =====================================================================

\maketitle

\begin{abstract}
Absolute stroke risk estimation is conventionally confined to clinical settings
because it requires measurements that individuals cannot obtain unaided. We present
a multimodal system that produces a calibrated five-year absolute stroke
probability from three inputs available to any smartphone user: a structured
questionnaire, a sixty-second fingertip photoplethysmogram acquired with the
device camera and torch, and optional photographs of medication packaging or
laboratory reports. Absolute risk is anchored to a published Cox
proportional-hazards model derived from a prospective cohort of 23{,}691 adults
with adjudicated stroke outcomes, extended with covariate-adjusted adjunct terms
for population-specific risk factors. Unanswered questions are propagated as
distributions rather than imputed, yielding an interval estimate together with a
variance-based ranking of which unmeasured input would most reduce uncertainty.
Rhythm assessment uses a classical interval-statistics classifier trained on the
MIT-BIH Atrial Fibrillation Database, achieving an area under the receiver
operating characteristic curve of 0.990 under leave-one-patient-out validation. We
introduce an arrhythmia-aware signal quality assessment, motivated by the
observation that the spectral and rate-stability metrics conventionally used to
gate photoplethysmographic heart-rate estimation are structurally incompatible
with arrhythmia detection. Independent evaluation on 300 smartphone recordings
with synchronous electrocardiographic ground truth yields 90\% beat-detection
sensitivity, 90\% positive predictive value, and 28.4\,ms inter-beat interval
error on expert-graded recordings. Evaluated at the deployed sixty-second window on
fingertip photoplethysmography from 35 adults with rhythm labels, the classifier attains
97.0\% sensitivity and 95.6\% specificity, statistically unchanged when the signal is
decimated to the 30\,Hz frame rate of a smartphone camera. Screening performance gives a
negative predictive value above 99.9\% at population prevalence, supporting an
asymmetric decision rule in which a regular rhythm resolves atrial fibrillation
within the risk model while an irregular rhythm produces a referral rather than a
classification.
\end{abstract}

\begin{IEEEkeywords}
stroke risk prediction, photoplethysmography, atrial fibrillation screening,
mobile health, survival analysis, uncertainty quantification, optical character
recognition
\end{IEEEkeywords}

%% =====================================================================
\section{Introduction}
%% =====================================================================

Stroke is among the leading causes of death and acquired disability worldwide, and
a substantial fraction of its burden is attributable to modifiable risk factors.
Established risk equations such as the Framingham Stroke Risk Profile translate a
set of clinical measurements into an absolute probability of incident stroke over
a defined horizon. Their practical reach is limited by the inputs they require:
systolic blood pressure measured by a trained observer, laboratory-confirmed
diabetes status, and electrocardiographically documented atrial fibrillation.

Two consequences follow. First, individuals outside regular clinical contact
cannot obtain an estimate at all. Second, where such tools are made available as
self-assessment instruments, unknown inputs are typically replaced by population
averages, and the resulting number is presented with a confidence that its inputs
do not support.

This work addresses both. We describe a system that produces a calibrated
five-year absolute stroke probability entirely on a consumer smartphone, combining
three modalities that require no hardware beyond the device itself, and that
represents what the user does not know as uncertainty rather than as an
assumption.

The principal contributions are:

\begin{enumerate}
\item An \emph{equation-anchored fusion architecture} in which absolute
      calibration derives from a published prospective-cohort survival model, and
      each additional modality either resolves a variable the model already
      contains or contributes a documented, covariate-adjusted adjunct term. Every
      coefficient remains traceable to a published source.
\item \emph{Uncertainty propagation with value-of-information ranking}. Unanswered
      questions are carried as prior distributions and propagated by Monte Carlo,
      producing an interval rather than a point estimate. First-order Sobol indices
      rank the unmeasured inputs by the share of output variance each explains,
      converting missing data into an actionable measurement recommendation.
\item \emph{Arrhythmia-aware signal quality assessment}. We show that spectral
      signal-to-noise ratio and inter-window rate consistency---the two metrics
      conventionally used to gate photoplethysmographic heart-rate estimation---are
      structurally unsuitable for gating arrhythmia analysis, and give a
      replacement based on band-aggregate and time-domain criteria.
\item An \emph{asymmetric decision rule} derived from the predictive-value
      asymmetry of low-prevalence screening, under which the rhythm module may
      exclude atrial fibrillation from the risk computation but never assert it.
\item \emph{Population-specific extension}, demonstrated by incorporating
      smokeless tobacco use---a dominant exposure in South Asia that Western-derived
      equations omit---as a separately weighted adjunct term.
\end{enumerate}

%% =====================================================================
\section{Related Work}
%% =====================================================================

\subsection{Absolute Risk Equations}
The Framingham Stroke Risk Profile \cite{wolf1991} established points-based
estimation of ten-year stroke probability from age, systolic blood pressure,
antihypertensive treatment, diabetes, smoking, prevalent cardiovascular disease,
atrial fibrillation and electrocardiographic left ventricular hypertrophy. A
revision \cite{dufouil2017} recalibrated the model against contemporary cohorts,
reporting concordance statistics between 0.62 and 0.78 across derivation and
external validation samples. Composite cardiovascular instruments such as the
Pooled Cohort Equations and QRISK3 include stroke within a broader endpoint.

\subsection{Machine Learning on Survey Data}
Gradient-boosted and ensemble classifiers trained on large cross-sectional health
surveys report discrimination exceeding 0.82 for stroke history. Such models
predict \emph{prevalent} self-reported stroke rather than incident events; because
the training data carry no time axis, their output is not a forward risk and is
not directly comparable to survival-model concordance.

\subsection{Smartphone Photoplethysmography}
Contact photoplethysmography using a fingertip placed over a smartphone camera
with the torch illuminated has been validated for atrial fibrillation screening in
several prospective studies, with reported sensitivity between 91.5\% and 96.3\%
and specificity between 97\% and 99.6\%
\cite{brasier2019,proesmans2019}. A meta-analysis of ten diagnostic accuracy
studies covering 3{,}852 participants found that, when screening asymptomatic
populations, positive predictive value falls to approximately 20--40\% while
negative predictive value approaches 100\% \cite{oconnor2020}. This asymmetry is
a direct consequence of low disease prevalence and motivates the decision rule
adopted here.

\subsection{Interval-Based Arrhythmia Detection}
Detection of atrial fibrillation from inter-beat intervals is well established
using statistical irregularity descriptors---root mean square of successive
differences, Shannon entropy, sample entropy \cite{richman2000} and Poincar\'e
descriptors. These operate on interval sequences and are therefore agnostic to
whether those intervals originate from electrocardiography or photoplethysmography,
which permits the classifier and the acquisition path to be validated separately.

%% =====================================================================
\section{System Architecture}
%% =====================================================================

The system executes entirely within the browser on the user's device. No
questionnaire response, video frame or photograph is transmitted. Fig.~1
summarises the flow.

\begin{figure}[t]
\centering
\footnotesize
\setlength{\fboxsep}{6pt}
\fbox{\begin{minipage}{0.88\columnwidth}
\centering
\textbf{Stage 1 — Questionnaire}\\
age, sex, systolic BP, medication, diabetes,\\
smoking, smokeless tobacco, prior CVD, AF\\[2pt]
$\downarrow$\\[2pt]
\textbf{Stage 2 — Fingertip PPG (60\,s)}\\
camera $+$ torch $\rightarrow$ quality gate $\rightarrow$\\
beat detection $\rightarrow$ RR intervals $\rightarrow$ classifier\\[2pt]
$\downarrow$\\[2pt]
\textbf{Stage 3 — Document capture (optional)}\\
on-device OCR $\rightarrow$ lexicon match $\rightarrow$\\
user confirmation\\[2pt]
$\downarrow$\\[2pt]
\textbf{Cox model $\rightarrow$ 5-year absolute risk}\\
point estimate or interval $+$ value of information
\end{minipage}}
\caption{Three-stage acquisition. Each stage may only resolve a variable left
unknown by the preceding stage; no stage overrides a direct user response.}
\end{figure}

A single invariant governs fusion: \emph{a later stage may resolve an unknown, but
never overrides an explicit user response, and never asserts a positive finding
that its evidence cannot support.}

%% =====================================================================
\section{Absolute Risk Estimation}
%% =====================================================================

\subsection{Model Form}
Five-year risk is computed from a Cox proportional-hazards model,
\begin{equation}
\hat{p}(t=5 \mid \mathbf{x}) = 1 - S_0(5)^{\exp\left(\mathrm{LP}(\mathbf{x}) - M\right)},
\end{equation}
where $\mathrm{LP}(\mathbf{x}) = \sum_j \beta_j x_j$, $\beta_j = \ln(\mathrm{HR}_j)$,
and $M = \sum_j \beta_j \bar{x}_j$ centres the predictor on the derivation cohort
mean. Diabetes enters through one of two age-conditioned coefficients, reflecting
the single retained interaction in the source model; systolic pressure enters
through a treated or untreated coefficient according to medication status.

\subsection{Coefficient Selection}
Two coefficient sets are published for the revised profile. We adopt the REGARDS
estimates as the operational default on four grounds: the derivation sample is
23{,}691 against 5{,}072; follow-up is exactly five years, matching the prediction
horizon without extrapolation; events were adjudicated under six-monthly active
surveillance; and the cohort includes substantial representation permitting
race-stratified baseline rates. Table~\ref{tab:coef} lists the hazard ratios.

\begin{table}[t]
\caption{Hazard ratios, REGARDS derivation sample, age $\geq$ 55}
\label{tab:coef}
\centering
\footnotesize
\begin{tabular}{lcc}
\toprule
Risk factor & Men & Women \\
\midrule
Age (per 10 years)          & 1.72 & 1.83 \\
Current smoking             & 2.00 & 1.77 \\
Prevalent CVD               & 1.31 & 1.52 \\
Atrial fibrillation         & 1.62 & 1.57 \\
Age $\geq$ 65               & 1.34 & 1.44 \\
Diabetes, age $<$ 65        & 1.94 & 1.66 \\
Diabetes, age $\geq$ 65     & 1.15 & 1.11 \\
Antihypertensive treatment  & 1.26 & 1.19 \\
SBP, untreated (per 10\,mmHg) & 1.09 & 1.16 \\
SBP, treated (per 10\,mmHg)   & 1.05 & 1.18 \\
\midrule
\textit{Adjunct term} & \multicolumn{2}{c}{} \\
Smokeless tobacco           & 1.35 & 1.35 \\
\bottomrule
\end{tabular}
\end{table}

\subsection{Recovery of Baseline Survival}
The source publication reports Kaplan--Meier cohort risk but not $S_0$. Because
risk is convex in the linear predictor, the mean of risk exceeds risk at the mean,
so equating $S_0$ with $1 - \mathrm{KM}$ introduces a systematic bias. We instead
construct a synthetic cohort matching the published covariate marginals and solve
numerically for the $S_0$ reproducing the observed mean risk:
\begin{equation}
S_0(5) : \quad \mathbb{E}_{\mathbf{x}}\!\left[1 - S_0(5)^{\exp(\mathrm{LP}(\mathbf{x})-M)}\right] = \mathrm{KM}(5).
\end{equation}
This yields $S_0(5) = 0.97310$ for men and $0.97803$ for women. The direct
substitution would give $0.969$ and $0.973$, overstating risk for a
cohort-representative individual by approximately 15\% in relative terms.
Reconstruction accuracy is verified by confirming that the mean predicted risk over
an independently seeded synthetic cohort reproduces the published five-year
Kaplan--Meier estimates of 0.031 and 0.027 to within 0.002.

\subsection{Adjunct Terms}
Risk factors absent from the derivation cohort are incorporated as additional
linear-predictor terms using published covariate-adjusted hazard ratios. Such
terms are centred at zero rather than at the target-population prevalence: the
derivation cohort is the reference against which $S_0$ was calibrated, and where
an exposure was negligible in that cohort, non-exposure already constitutes the
reference level. Centring on target prevalence would apply an unwarranted
reduction to unexposed individuals. Cohort calibration is verified to be
unaffected.

We apply this to smokeless tobacco, with a pooled stroke relative risk of 1.35
(95\% CI 1.20--1.50) for chewing-tobacco users \cite{gupta2019}. The same
meta-analysis reports the overall association to be essentially unchanged under
strict adjustment for smoking (1.17 unadjusted versus 1.18 adjusted), supporting
treatment of this exposure as separable from combustible tobacco. It is queried
independently because its effect is materially smaller than that of smoking
(HR 2.00 in men); combining the two would overstate risk for exclusive smokeless
users by approximately 50\%.

\subsection{Uncertainty Propagation}
Where a response is unknown, the variable is sampled from a population prior and
propagated by Monte Carlo, and the system reports the 10th--90th percentile
interval. A fully specified profile collapses the interval to a point.

To convert missing data into guidance, we compute the first-order Sobol index for
each unknown input,
\begin{equation}
S_j = \frac{\mathrm{Var}(R) - \mathbb{E}_{x_j}\!\left[\mathrm{Var}(R \mid X_j = x_j)\right]}{\mathrm{Var}(R)},
\end{equation}
the share of output variance attributable to that input. Variance-based
decomposition is bounded in $[0,1]$ and is well defined for the non-additive
response surface, and the resulting ranking is subject-specific rather than fixed.

\subsection{Counterfactual Guidance}
Modifiable factors are reported with the risk each change would remove. Changes in
medication \emph{status} are deliberately excluded from this set. At fixed measured
systolic pressure the fitted model assigns treated individuals a lower hazard than
untreated individuals, reflecting the flatter observed pressure--risk gradient in
treated cohorts rather than a causal treatment effect; presenting it as an
actionable lever would misrepresent the evidence. The pressure value itself is
offered instead.

%% =====================================================================
\section{Rhythm Assessment from Smartphone Photoplethysmography}
%% =====================================================================

\subsection{Acquisition}
The user places a fingertip over the rear camera and torch for sixty seconds.
Illuminated contact photoplethysmography is used exclusively; ambient-light facial
measurement is not offered, as it depends on skin reflectance and is markedly more
sensitive to motion and illumination change. Per-frame mean intensity is computed
over a central region of interest for each colour channel, and the channel
carrying the cleanest pulse is selected automatically.

\subsection{Sub-Sample Beat Localisation}
\label{sec:subsample}
At a typical capture rate of 30\,frames\,s$^{-1}$, beat timing is quantised to
33\,ms---comparable in magnitude to the beat-to-beat variability that distinguishes
sinus rhythm from atrial fibrillation. Fitting a parabola through each candidate
peak and its immediate neighbours recovers the extremum to a fraction of a frame.
On a synthetic signal whose period is incommensurate with the frame interval, this
reduces interval error from 15.9\,ms to 0.04\,ms, rendering frame rate a
non-limiting factor.

\subsection{Arrhythmia-Aware Quality Assessment}
\label{sec:quality}
Photoplethysmographic quality gating conventionally relies on two families of
metric: spectral signal-to-noise ratio, and consistency of the estimated rate
across sliding windows. Both are appropriate when the objective is a heart rate.
Both are structurally unsuitable when the objective is rhythm classification:

\begin{itemize}
\item \textbf{Inter-window rate consistency} penalises a rate that varies between
      windows. In atrial fibrillation the rate genuinely varies between windows;
      the criterion therefore rejects the target condition.
\item \textbf{Spectral signal-to-noise ratio} quantifies the sharpness of the
      dominant spectral peak. Irregular intervals distribute pulse energy across
      the band, so atrial fibrillation has no sharp peak by construction. A clean
      synthetic fibrillatory waveform measures a negative signal-to-noise ratio
      under this definition.
\end{itemize}

Both metrics conflate the signal of interest with noise. We therefore gate rhythm
analysis on criteria that are band-aggregate or purely temporal, and which respond
to motion without responding to rhythm:

\begin{enumerate}
\item \textbf{In-band power fraction}: the proportion of detrended signal variance
      falling within the pulse band, which detects the presence of a pulse at
      \emph{any} rate within the band rather than at one frequency.
\item \textbf{Photometric stability}: relative variation in mean frame intensity,
      which responds directly to finger displacement and illumination change.
\item \textbf{Beat amplitude consistency}: the coefficient of variation of detected
      peak amplitudes. Motion artefacts produce heterogeneous peak morphology;
      atrial fibrillation alters interval spacing while leaving pulse morphology
      comparatively uniform.
\item \textbf{Interval plausibility rate}: the proportion of detected intervals
      falling outside physiological bounds, distinguishing failed beat detection
      from genuine irregularity within plausible bounds.
\end{enumerate}

\subsection{Features and Classifier}
Eight scale-free descriptors are computed over each analysis window: normalised
root mean square of successive differences, coefficient of variation, the
proportion of successive differences exceeding 50\,ms and 70\,ms, normalised
Shannon entropy of the interval histogram, sample entropy, the Poincar\'e
$\mathrm{SD1}/\mathrm{SD2}$ ratio, and mean rate. A standardised logistic
regression is fitted over these features. A linear model is retained deliberately:
the decision is auditable, the exported parameterisation is 4\,kB, and inference
requires no runtime dependency.

Sample entropy is undefined when no template pairs match within tolerance, a
condition arising at both extremes of regularity---highly irregular sequences, where
no matches exist, and near-constant sequences, where the matching tolerance
approaches zero. We substitute the conventional upper bound
$-\ln\!\left(2/[(N{-}m{-}1)(N{-}m)]\right)$, mapping both limiting cases to maximal
entropy and admitting the full range of rhythm into training.

\subsection{Asymmetric Decision Rule}
\label{sec:asym}
Because predictive value is prevalence-dependent, a screening instrument with
high sensitivity and specificity may still yield a positive predictive value far
below certainty. The module therefore emits one of three outcomes:

\begin{itemize}
\item \textbf{Regular}: atrial fibrillation is resolved to negative within the risk
      model, and the interval estimate narrows accordingly.
\item \textbf{Irregular}: the risk estimate is left unchanged and the user is
      advised to seek electrocardiography. No classification is asserted.
\item \textbf{Inconclusive}: quality criteria are not met; no inference is drawn.
\end{itemize}

This is not a conservatism heuristic but a direct consequence of the predictive
values reported in Section~\ref{sec:results}.

%% =====================================================================
\section{Medical Document Analysis}
%% =====================================================================

The optional third modality accepts a photograph of medication packaging, a
prescription, or a laboratory report. Optical character recognition executes
on-device, preserving the property that no image leaves the handset.

Recognised text is matched against a curated lexicon of 56 generic agents and 191
proprietary names, weighted toward South Asian markets where proprietary labelling
predominates. Matching is tolerant to the substitution errors characteristic of
curved, specular packaging surfaces, using normalised tokens compared under a
length-dependent bounded edit distance. Laboratory values are extracted by
proximity of a numeric literal to a recognised assay name, with range validation
to reject spurious matches.

Each agent class maps to a model variable with an associated inference strength.
Classes with a single dominant indication propose the corresponding variable
directly. Classes with multiple indications---anticoagulants, which may be
prescribed for fibrillation, venous thrombosis or prosthetic valves---raise a
clarifying question rather than a value. Corroborating evidence from independent
sources is accumulated into a single proposal.

No extracted value modifies the risk computation without explicit user
confirmation. This constraint reduces recognition error from a correctness concern
to an interaction concern: a misread becomes a question the user declines.

%% =====================================================================
\section{Evaluation Methodology}
%% =====================================================================

\subsection{Rhythm Classifier}
The classifier was trained and evaluated on the MIT-BIH Atrial Fibrillation
Database, comprising 25 long-term recordings with verified beat annotations and
rhythm episode labels \cite{goldberger2000}. Intervals were segmented into
60-second windows at 30-second
hops; a window was labelled fibrillatory when at least 90\% of its beats fell
within an annotated episode, yielding 29{,}018 windows of which 37.2\% were
fibrillatory. Validation was leave-one-record-out, so no patient contributed to
both training and evaluation.

\subsection{Degradation Analysis}
To characterise tolerance to acquisition error, features were recomputed from the
annotated beat times after injecting timing perturbation and beat omission, and the
full leave-one-record-out procedure was repeated at each level.

\subsection{Acquisition Path}
The acquisition path was evaluated independently on the Brno University of
Technology Smartphone PPG Database, comprising fingertip recordings captured on
consumer handsets at 30\,frames\,s$^{-1}$ with synchronous 1\,kHz
electrocardiography and verified QRS annotations \cite{nemcova2021}. Analysis was
restricted to fingertip acquisitions, yielding 300 recordings across eight
annotated activity conditions.

Beat-detection accuracy was assessed by temporal correspondence rather than
sequence subtraction: the pulse transit delay was estimated per recording, each
detected peak matched to an electrocardiographic beat within a 100\,ms tolerance,
and interval error computed only across pairs whose endpoints both matched and
were consecutive in both streams. This separates beat-detection sensitivity and
positive predictive value from interval precision. The matching procedure was
itself verified against synthetic sequences with known transit delay, known
timing jitter, and deliberately omitted and inserted beats.

\subsection{Deployed Window Length}
The fingertip acquisition corpus provides ten-second clips, whereas the classifier
operates over sixty seconds. Evaluation at the deployed configuration used an
independent smartphone corpus recorded at the Gda\'{n}sk University of Technology
\cite{bancerewicz2025}, comprising one sixty-one-second acquisition with aligned
electrocardiography and three continuous recordings of approximately ten minutes from
healthy volunteers, segmented into non-overlapping sixty-second windows.

\subsection{Rhythm Classification at the Deployed Configuration}
\label{sec:mimicmethod}
The MIMIC PERform AF corpus \cite{charlton2022} provides simultaneous fingertip
photoplethysmography and electrocardiography from 35 critically ill adults, 19 in atrial
fibrillation and 16 not, over twenty minutes each at 125\,Hz. It is the only openly
accessible rhythm-labelled corpus that is both fingertip---transmission through the
digit, the optical geometry of a torch-illuminated camera capture---and long enough to
segment at the deployed sixty-second window. Recordings were divided into
non-overlapping sixty-second windows, discarding windows in which the sensor was
disconnected, yielding 659 windows. To separate the contribution of frame rate from that
of the sensing site, the evaluation was repeated on the same windows decimated to
30\,Hz, the rate at which the system deploys. Rhythm labels are assigned per recording,
so results are additionally aggregated per subject.

\subsection{Rhythm Classification on Wrist Photoplethysmographic Intervals}
Classifier sensitivity established on MIT-BIH is measured on intervals derived from the
QRS complex. To measure it on intervals derived from a pulse wave instead, the frozen
classifier was applied to the held-out test split of the DeepBeat corpus
\cite{torressoto2020}: 17{,}617 wrist photoplethysmographic segments of twenty-five
seconds at 32\,Hz drawn from 22 subjects, each segment carrying a rhythm label and a
three-level signal-quality label. The rhythm label is constant within a subject, so the
effective sample size is the subject count and results are additionally aggregated per
subject. The training and validation splits were not accessed.

\subsection{Implementation Verification}
The deployed implementation is JavaScript; the reference implementation is Python.
Agreement is enforced by replaying 10{,}764 reference-computed risk values and 358
reference-computed rhythm features through the deployed code, with a tolerance of
$10^{-12}$. Observed worst-case divergence was $1.1 \times 10^{-16}$. Model
parameters were fixed and cryptographically hashed prior to acquisition-path
evaluation, and verified unchanged afterwards.

%% =====================================================================
\section{Results}
\label{sec:results}
%% =====================================================================

\subsection{Rhythm Classification}
Under leave-one-record-out validation the classifier achieved an area under the
receiver operating characteristic curve of 0.9903. At an operating point selected
for 95\% sensitivity, specificity was 97.4\%. Per-record discrimination ranged from
0.931 to 1.000.

\begin{table}[t]
\caption{Classifier tolerance to acquisition timing error}
\label{tab:noise}
\centering
\footnotesize
\begin{tabular}{lccc}
\toprule
Condition & AUROC & Sens. & Spec. \\
\midrule
Annotated beat times (reference)      & 0.9903 & 95\% & 97.4\% \\
30\,fps quantisation, no interpolation & 0.9918 & 95\% & 96.9\% \\
30\,fps with sub-sample interpolation  & 0.9900 & 95\% & 96.9\% \\
60\,fps, 5\,ms jitter                  & 0.9906 & 95\% & 97.3\% \\
15\,ms jitter                          & 0.9880 & 95\% & 95.7\% \\
25\,ms jitter                          & 0.9826 & 95\% & 92.5\% \\
40\,ms jitter                          & 0.9699 & 95\% & 87.4\% \\
8\,ms jitter, 5\% beats omitted        & 0.9900 & 95\% & 96.7\% \\
\bottomrule
\end{tabular}
\end{table}

Table~\ref{tab:noise} shows discrimination is preserved to approximately 25\,ms of
per-beat timing error, and that frame-rate quantisation alone is not limiting.
Random beat omission at 5\% has negligible effect, indicating that the
interval-statistic descriptors are robust to sparse detection failure.

\subsection{Acquisition Path}

\begin{table}[t]
\caption{Beat detection and interval accuracy, smartphone acquisition\\
against synchronous electrocardiography ($n = 300$)}
\label{tab:butppg}
\centering
\footnotesize
\begin{tabular}{lcccc}
\toprule
Subset & Sens. & PPV & RR err. & HR err. \\
\midrule
Expert-graded good ($n{=}100$) & 90\% & 90\% & 28.4\,ms & 1.0\,bpm \\
Accepted by quality gate       & 85\% & 92\% & 28.9\,ms & 2.1\,bpm \\
Expert-graded poor ($n{=}200$) & 40\% & 57\% & 63.2\,ms & 19.7\,bpm \\
Rejected by quality gate       & 44\% & 63\% & 47.6\,ms & 15.4\,bpm \\
\bottomrule
\end{tabular}
\end{table}

Table~\ref{tab:butppg} reports performance stratified by the database's own expert
quality annotations and by the system's quality gate. On expert-graded good
recordings, beat-detection sensitivity and positive predictive value both reach
90\%, with interval error of 28.4\,ms and heart-rate error of 1.0\,bpm; 83\% of
accepted recordings fall within the 5\,bpm tolerance of IEC 60601-2-27.

The quality gate discriminates concordantly with the expert annotations,
accepting recordings with more than twice the beat-detection sensitivity of those
it rejects (85\% versus 44\%) and admitting expert-graded good recordings at
approximately twice the rate of expert-graded poor recordings.

\subsection{Deployed Window Length}
At the deployed sixty-second configuration, specificity on continuous recordings from
healthy volunteers was 91.7\%---22 of 24 windows called regular. This is a direct
measurement at the operating configuration rather than an extrapolation from shorter
clips, and it agrees to within 0.7 percentage points with the 92.4\% projected below
from acquisition error measured on an unrelated corpus.

\subsection{Rhythm Classification at the Deployed Configuration}

\begin{table}[t]
\caption{Frozen classifier on fingertip photoplethysmography at the\\
deployed sixty-second window ($n = 659$ windows, 35 subjects)}
\label{tab:mimic}
\centering
\footnotesize
\begin{tabular}{lccc}
\toprule
Sampling rate & Sensitivity & Specificity & AUROC \\
\midrule
30\,Hz (deployed frame rate) & 97.0\% & 95.6\% & 0.9918 \\
125\,Hz (native)             & 96.7\% & 94.9\% & 0.9902 \\
\midrule
\emph{MIT-BIH, ECG intervals} & \emph{95\%} & \emph{97.4\%} & \emph{0.9903} \\
\bottomrule
\end{tabular}
\end{table}

Table~\ref{tab:mimic} evaluates the frozen classifier at both the deployed window length
and the deployed anatomical site. Sensitivity is 97.0\% (95\% CI 94.7--98.3\%) and
specificity 95.6\% (92.6--97.4\%) at 30\,Hz, with an area under the curve of 0.9918
against 0.9903 on electrocardiographic intervals. Per subject, all 19 fibrillatory
subjects and 15 of 16 sinus-rhythm subjects were classified correctly. Every window was
scored; none was rejected as unusable.

Decimating from 125\,Hz to 30\,Hz changes sensitivity by 0.3 points and specificity by
0.7---within the confidence intervals of both. Camera frame rate is therefore not the
limiting factor in interval-based rhythm discrimination, which is the outcome the
sub-sample interpolation of Section~\ref{sec:subsample} is intended to produce, measured
here rather than assumed.

\subsection{Rhythm Classification on Wrist Photoplethysmographic Intervals}

\begin{table}[t]
\caption{Frozen classifier on wrist photoplethysmography, stratified\\
by the corpus signal-quality label ($n = 17{,}614$ scored)}
\label{tab:deepbeat}
\centering
\footnotesize
\begin{tabular}{lccc}
\toprule
Signal quality & Segments & Sensitivity & Specificity \\
\midrule
Excellent  & 3{,}246  & 91.5\% & 92.4\% \\
Acceptable & 2{,}032  & 91.7\% & 81.0\% \\
Poor       & 12{,}336 & 76.1\% & 40.6\% \\
\bottomrule
\end{tabular}
\end{table}

Table~\ref{tab:deepbeat} reports the frozen classifier applied to intervals derived from
a pulse wave rather than from electrocardiography. On excellent-quality segments
sensitivity is 91.5\% and specificity 92.4\%, against 95\% and 97.4\% on
electrocardiographic intervals: the transition from QRS timing to pulse-wave timing,
with its softer peaks and respiratory amplitude modulation, costs approximately three
points of sensitivity and five of specificity. Performance declines monotonically with
signal quality, and 70\% of this corpus is graded poor, reflecting continuous wrist wear
rather than deliberate fingertip acquisition.

Aggregated per subject over excellent-quality segments, all four fibrillatory subjects
and all six sinus-rhythm subjects were classified correctly. With ten subjects these are
reported as counts rather than as rates.

\subsection{Projected Screening Performance}
The measured interval error of 28.4\,ms corresponds, under a Gaussian error model,
to a per-beat timing dispersion of approximately 25\,ms. Interpolating
Table~\ref{tab:noise} at that operating point projects a specificity of 92.4\% at
95\% sensitivity.

\begin{table}[t]
\caption{Projected predictive values by screening prevalence}
\label{tab:ppv}
\centering
\footnotesize
\begin{tabular}{lcc}
\toprule
Population (AF prevalence) & PPV & NPV \\
\midrule
General adults, age $\geq$ 55 \quad (2\%)  & 20.4\% & 99.89\% \\
Adults, age $\geq$ 65 \quad (5\%)          & 39.8\% & 99.72\% \\
Higher-risk, age $\geq$ 75 \quad (10\%)    & 58.3\% & 99.40\% \\
\bottomrule
\end{tabular}
\end{table}

Table~\ref{tab:ppv} gives the resulting predictive values. These fall within the
20--40\% positive and near-100\% negative predictive values reported for validated
smartphone screening applications \cite{oconnor2020}, despite being derived
independently---from classifier performance on electrocardiographic data combined
with acquisition error measured on a separate smartphone corpus. The agreement
between two unrelated derivations supports the validity of the projection.

The photoplethysmographic operating point of Table~\ref{tab:deepbeat} provides a third,
independent route to the same quantity: 91.5\% sensitivity at 92.4\% specificity yields
19.7\% positive and 99.81\% negative predictive value at 2\% prevalence, within one
percentage point of Table~\ref{tab:ppv} despite deriving from a different corpus, a
different anatomical site and a different sensing modality.

Table~\ref{tab:ppv} is nonetheless conservative. Applying the measured fingertip
operating point of Table~\ref{tab:mimic} in place of the projection gives 31.0\%
positive and 99.94\% negative predictive value at 2\% prevalence. The qualitative
asymmetry that governs the decision rule is unchanged: a negative result remains
dependable and a positive result remains, more often than not, wrong.

These values establish the decision rule of Section~\ref{sec:asym} quantitatively.
A negative predictive value of 99.9\% renders exclusion of atrial fibrillation
dependable; a positive predictive value of 20\% at population prevalence renders
assertion unsupportable. The module is therefore permitted to exclude and required
to refer.

\subsection{Risk Model Behaviour}
Reconstruction reproduces the published five-year cohort risks to within 0.002 for
both sexes and both coefficient sets. Risk is monotone in age, systolic pressure
and each binary factor. For a 68-year-old man with no other information supplied,
the interval estimate spans 2.2--5.0\%, and the variance decomposition attributes
52\% of that uncertainty to smoking status and 10\% to systolic pressure; for a
58-year-old woman the ranking inverts, with systolic pressure accounting for 36\%,
reflecting the steeper pressure gradient in the female coefficient set. The
recommendation is thus subject-specific rather than a fixed ordering.

%% =====================================================================
\section{Discussion}
%% =====================================================================

The central design question for a self-administered risk instrument is not which
additional sensor to add, but how to behave when required inputs are unavailable.
Substituting population means for unknown inputs produces a confident number from
fabricated data. Representing them as distributions produces a wider but honest
estimate, and---through variance decomposition---an actionable statement of which
single measurement would most reduce that width. This reframes missing data from a
defect of the instrument into its primary output.

The quality-gating result of Section~\ref{sec:quality} generalises beyond this
system. Quality metrics are commonly transferred between photoplethysmographic
tasks on the assumption that a clean signal is clean for all purposes. Where the
target phenomenon is itself a departure from regularity, metrics defined in terms
of regularity---spectral peak sharpness, rate stability---cannot separate signal from
artefact. Any pipeline gating arrhythmia analysis with heart-rate quality criteria
will preferentially discard its positive cases, and will do so silently.

The predictive-value asymmetry has a design consequence frequently left implicit.
A screening instrument reporting 95\% sensitivity and 92\% specificity is
legitimately described as accurate, yet at 2\% prevalence four of five positive
indications are false. Systems that report such a result as a finding will
mislead the majority of those they flag. Reporting it as a referral preserves the
clinical value of the true positives while making no claim the evidence cannot
support.

Finally, population-specific extension is achievable without abandoning a validated
equation. Adding smokeless tobacco as a zero-centred adjunct term with a published
adjusted hazard ratio leaves the derivation-cohort calibration intact while
capturing an exposure affecting 29.6\% of men and 12.8\% of women in the target
population \cite{gats2017}, which the underlying Western-derived model does not represent at all.

%% =====================================================================
\section{Limitations and Future Work}
%% =====================================================================

\textbf{Sensitivity and specificity derive from distinct corpora.} Classifier
sensitivity is established on electrocardiographic intervals and corroborated on both
fingertip and wrist photoplethysmographic intervals; acquisition fidelity and specificity are established on
smartphone fingertip recordings. No public corpus provides smartphone \emph{fingertip}
photoplethysmography with adjudicated fibrillation labels, so the end-to-end performance
of Table~\ref{tab:ppv} remains assembled from separate measurements rather than observed
directly. A prospective acquisition with synchronous electrocardiography would resolve
this and is the principal item of future work.

\textbf{Dependence on signal quality.} Table~\ref{tab:deepbeat} shows specificity falling
to 40.6\% on segments the corpus grades poor. Reported performance is therefore
conditional on the quality gate admitting only adequate recordings. That gate was
characterised on deliberate fingertip acquisition, and its behaviour on continuous wrist
wear---the modality of Table~\ref{tab:deepbeat}---has not been separately quantified.

\textbf{Analysis window.} Classification has been measured at the deployed sixty-second
window on fingertip photoplethysmography (Table~\ref{tab:mimic}) and on continuous
smartphone recordings from healthy volunteers, but the smartphone \emph{camera}
acquisition corpus provides only ten-second clips. The step that remains inferred rather
than observed is therefore the camera itself, not the window length or the sensing
site.

\textbf{Cohort representativeness.} The acquisition corpus comprises fifty
volunteers with a mean age of 33 years, and does not report Fitzpatrick skin
phototype. Since melanin attenuates the optical signal, performance across skin
tone remains uncharacterised and warrants explicit stratification in prospective
evaluation.

\textbf{Derivation range and geography.} The risk equation is derived on adults
aged 55--84 in United States cohorts; the system labels estimates outside that
range as extrapolation. Relative effects transport between populations more
reliably than absolute baselines, and recalibration of $S_0$ to age- and
sex-specific local incidence is a data requirement rather than a modelling one.
Reported calibration for this model family is weaker in Black participants than in
white participants, and the system discloses this in its output.

\textbf{Adjunct term independence.} Adjunct hazard ratios are drawn from
meta-analyses with heterogeneous covariate adjustment, and are combined
multiplicatively with the base model under an independence assumption that is
defensible where the exposure is absent from the derivation cohort, but is not
directly verified.

\textbf{Regulatory scope.} Presentation of an absolute stroke probability to a lay
user constitutes a medical device claim in most jurisdictions. The system is
presented as a research instrument and is not offered for clinical use.

%% =====================================================================
\section{Conclusion}
%% =====================================================================

We have presented a multimodal stroke risk system producing a calibrated five-year
absolute probability from a questionnaire, a sixty-second fingertip
photoplethysmogram, and optional medical document capture, executing entirely on a
consumer smartphone with no data transmission. Absolute calibration is anchored to
a published prospective-cohort survival model with fully traceable coefficients;
unknown inputs are propagated as uncertainty and ranked by variance contribution
rather than imputed; and rhythm assessment employs quality criteria constructed
specifically for arrhythmia rather than inherited from heart-rate estimation.

Rhythm classification achieves an area under the curve of 0.990 under
leave-one-patient-out validation on electrocardiographic intervals, and 0.992 at the
deployed sixty-second window on fingertip photoplethysmography from an independent
corpus, where sensitivity and specificity are 97.0\% and 95.6\%. The acquisition path
achieves 90\% beat-detection sensitivity with 28.4\,ms interval error on expert-graded
smartphone recordings. That the fingertip result is unchanged by decimation to 30\,Hz
indicates camera frame rate is not the limiting factor in interval-based rhythm
discrimination. The resulting predictive-value asymmetry---above 99.9\% negative against
20--31\% positive at population prevalence---directly determines a decision rule under
which the system may exclude atrial fibrillation from the risk computation but
refers rather than asserts it. We argue this pattern, in which the evidentiary
strength of a screening result governs the action it is permitted to take, is the
appropriate organising principle for consumer-facing diagnostic inference.

%% =====================================================================
\begin{thebibliography}{00}

\bibitem{wolf1991}
P.~A. Wolf, R.~B. D'Agostino, A.~J. Belanger, and W.~B. Kannel,
``Probability of stroke: a risk profile from the Framingham Study,''
\emph{Stroke}, vol.~22, no.~3, pp.~312--318, 1991.

\bibitem{dufouil2017}
C.~Dufouil, A.~Beiser, L.~A. McLure, P.~A. Wolf, C.~Tzourio, V.~J. Howard,
et~al.,
``Revised Framingham Stroke Risk Profile to reflect temporal trends,''
\emph{Circulation}, vol.~135, no.~12, pp.~1145--1159, 2017.

\bibitem{brasier2019}
N.~Brasier, C.~J. Raichle, M.~Dörr, A.~Becke, V.~Nohturfft, S.~Weber, et~al.,
``Detection of atrial fibrillation with a smartphone camera: first prospective,
international, two-centre, clinical validation study (DETECT AF PRO),''
\emph{EP Europace}, vol.~21, no.~1, pp.~41--47, 2019.

\bibitem{proesmans2019}
T.~Proesmans, C.~Mortelmans, R.~Van Haelst, F.~Verbrugge, P.~Vandervoort, and
B.~Vaes,
``Mobile phone-based use of the photoplethysmography technique to detect atrial
fibrillation in primary care: diagnostic accuracy study of the FibriCheck app,''
\emph{JMIR mHealth and uHealth}, vol.~7, no.~3, e12284, 2019.

\bibitem{oconnor2020}
J.~W. O'Sullivan, S.~Grigg, W.~Crawford, M.~P. Turakhia, M.~Perez,
E.~Ingelsson, M.~T. Wheeler, and J.~P.~A. Ioannidis,
``Accuracy of smartphone camera applications for detecting atrial fibrillation: a
systematic review and meta-analysis,''
\emph{JAMA Network Open}, vol.~3, no.~4, e202064, 2020.

\bibitem{richman2000}
J.~S. Richman and J.~R. Moorman,
``Physiological time-series analysis using approximate entropy and sample
entropy,''
\emph{American Journal of Physiology---Heart and Circulatory Physiology},
vol.~278, no.~6, pp.~H2039--H2049, 2000.

\bibitem{gupta2019}
R.~Gupta, S.~Gupta, S.~Sharma, D.~N. Sinha, and R.~Mehrotra,
``Association of smokeless tobacco and cerebrovascular accident: a systematic
review and meta-analysis of global data,''
\emph{Journal of Public Health}, vol.~42, no.~2, pp.~e150--e157, 2020.
[Corrigendum: \emph{J. Public Health}, vol.~43, no.~1, p.~e115, 2021.]

\bibitem{nemcova2021}
A.~N\v{e}mcov\'a, E.~Vargov\'a, R.~Smisek, L.~Mar\v{s}\'anov\'a,
L.~Smital, and M.~Vitek,
``Brno University of Technology Smartphone PPG Database (BUT PPG): annotated
dataset for PPG quality assessment and heart rate estimation,''
\emph{BioMed Research International}, vol.~2021, art.~3453007, 2021.

\bibitem{goldberger2000}
A.~L. Goldberger, L.~A.~N. Amaral, L.~Glass, J.~M. Hausdorff, P.~C. Ivanov,
R.~G. Mark, et~al.,
``PhysioBank, PhysioToolkit, and PhysioNet: components of a new research resource
for complex physiologic signals,''
\emph{Circulation}, vol.~101, no.~23, pp.~e215--e220, 2000.

\bibitem{charlton2022}
P. H. Charlton, K. Kotzen, E. Mej\'{i}a-Mej\'{i}a, P. J. Aston, K. Budidha, J. Mant,
C. Pettit, J. A. Behar, and P. A. Kyriacou, ``Detecting beats in the photoplethysmogram:
benchmarking open-source algorithms,'' \emph{Physiological Measurement}, vol. 43, art.
085007, 2022.

\bibitem{torressoto2020}
J. Torres-Soto and E. A. Ashley, ``Multi-task deep learning for cardiac rhythm detection
in wearable devices,'' \emph{npj Digital Medicine}, vol. 3, art. 116, 2020.

\bibitem{bancerewicz2025}
J. K. Bancerewicz, J. J. Kot\l{}owski, O. Lozovyy, J. B. Morawska, and M. Rz\k{e}sa,
``PPGbetter: smartphone photoplethysmography acquisition for heart-rate variability
research,'' software repository and accompanying recordings, Gda\'{n}sk University of
Technology, 2025.

\bibitem{gats2017}
Ministry of Health and Family Welfare, Government of India,
``Global Adult Tobacco Survey: India 2016--17 report,'' 2017.

\end{thebibliography}

\end{document}
